\documentclass[a4paper,12pt]{article}
\usepackage{color}
\usepackage{graphicx, gensymb}
\usepackage{amsfonts, cancel, amssymb}
\usepackage{amsmath, amsthm, array, esvect, esint}
\usepackage{typearea}
\usepackage[a4paper,left=2cm,right=2cm,top=2cm,bottom=3cm,bindingoffset=5mm]{geometry}
\newcommand\numberthis{\addtocounter{equation}{1}\tag{\theequation}}
\newcommand{\bra}{}
\newcommand{\kom}{}
\newcommand{\brak}{}
\newcommand{\braket}{}
\newcommand{\intx}{}
\newcommand{\intr}{}
\newcommand{\kla}{}
\newcommand{\klam}{}
\newcommand{\klammer}{}
\newcommand{\oper}{}
\newcommand{\tecxt}{}
\newcommand{\Bra}{}
\newcommand{\Ket}{}

\begin{document}

\title{5 Freie Energie}
\author{Joachim Schwardt}
\date{\today}
\maketitle

\section*{Aufgabe 5.1: Freie Energie und Barometrische Höhenformel}
Es gilt $F=U-TS$. Man kann diese über
\begin{align*}
F[\rho]&=\intx\int_{0}^{\infty}-\rho(z)mgz-k_BT(z)\rho(z)\log\rho(z)\,\mathrm{d}z
\end{align*}
berechnen. Die minimierende Verteilung erhält man durch ableiten:
\begin{align*}
0&=-mgz-k_BT\log\rho -k_BT&\Rightarrow &&\rho(z)&=\rho_0\exp\kla\left( {-\frac{mgz}{k_BT(z)}} \right)
\end{align*}
Dabei wurde $\rho$ als die Anzahl an Teilchen pro Höhe interpretiert. Allerdings ist die potentielle Energie genauer durch $V=-\frac{GMm}{(z+R_E)^2}$ gegeben.
\section*{Aufgabe 5.2: Polymer-Kette und entropische Kräfte}
Sei $N$ die Anzahl an Gliedern und $N_\pm$ die Anzahl in positive bzw. negative $x$-Richtung. Dann endet das letzte Glied der Kette auf $k=N_+-N_-$, wobei die Koordinate durch $x=k\cdot a$ gegeben ist. Bekannt ist aber auch $N=N_++N_-$, woraus direkt $N_+=\frac{N+k}{2}$ folgt. Die Wahrscheinlichkeit, dass das $N$-te Glied des Polymers auf der Koordinate $x=k\cdot a$ ($k\in\mathbb{Z}$) endet, beträgt somit
\begin{align*}
P(k;N)=\frac{1}{2^N}\binom{N}{N_+}=\frac{1}{2^N}\frac{N!}{(\frac{N+k}{2})!(\frac{N-k}{2})!}
\end{align*}
Dieser Term kann mit $N!\sim\sqrt{2\pi N}(N/e)^N$ folgendermaßen genähert werden:
\begin{align*}
P(k;N)&\sim\frac{1}{2^N}\frac{\sqrt{2\pi N}N^N}{\sqrt{2\pi(\frac{N+k}{2})}(\frac{N+k}{2})^{\frac{N+k}{2}}\sqrt{2\pi (\frac{N-k}{2})}(\frac{N-k}{2})^{\frac{N-k}{2}}}=\\
&=\frac{(1+k/N)^{-N/2}(1-k/N)^{-N/2}}{\sqrt{\frac{\pi}{2}(1+k/N)(1-k/N)}(\frac{1+k/N}{1-k/N})^{k/2}}
\end{align*}
Für den Grenzübergang $N\rightarrow\infty$ betrachtet man $\log P$ mit $\log(1+x)\sim x$ für $x\ll 1$:
\begin{align*}
\log P&\sim -\frac{N}{2}\log\kla\left( {1-\frac{k^2}{N^2}} \right)-\frac{k}{2}\log \kla\left( {1+\frac{k}{N}} \right)+\frac{k}{2}\log \kla\left( {1-\frac{k}{N}} \right)\kom\overset{\substack{{k\ll N} \\ |}}{\sim}\\
&\sim\frac{k^2}{2N}-\frac{k^2}{2N}-\frac{k^2}{2N}=-\frac{k^2}{2N}
\end{align*}
Daraus folgt eine kontinuierliche Näherung für große $N$ und $k\ll N$:
\begin{align*}
P(k;N)&\sim Ce^{-\frac{k^2}{2N}}=\frac{1}{\sqrt{2\pi N}}e^{-\frac{k^2}{2N}}
\end{align*}
Die Konstante erhält man dabei aus der Normierung.\\
Oder: $r=\sum_ja_j$ mit $a_j=a\cdot(2X_j-1)$ und $X_j\in\{0,1\}$ Bernoulli-verteilt mit $p=\frac{1}{2}$. Dann ist
\begin{align*}
r&=2a\sum_jX_j-Na
\end{align*}
Binomial-verteilt. Daraus folgt direkt
\begin{align*}
\bra\langle {r}\rangle&=2a\bra\langle {\sum_jX_j}\rangle-Na=2a\cdot\frac{N}{2}-Na=0\\
\bra\langle {r^2}\rangle&=4a^2\bra\langle {\kla\left( {\sum_jX_j} \right)^2}\rangle=4a^2\cdot N\cdot\frac{1}{2}\cdot\frac{1}{2}=Na^2=:r_0^2
\end{align*}
Für $r_0=$const geht die Verteilung bei $N\rightarrow\infty$ in eine Gauss-Verteilung über:
\begin{align*}
\rho(r)&\sim\frac{1}{\sqrt{2\pi r_0^2}}e^{-\frac{r^2}{2r_0^2}}
\end{align*}
Unter Vernachlässigung aller Wechselwirkungen ist $U=0$, und damit folgt aus $\Omega(r)=\rho(r)\Omega$ mit $\Omega=2^N$ die freie Energie zu
\begin{align*}
F&=U-TS=k_BT\sum_\Omega\frac{1}{\Omega(r)}\log\frac{1}{\Omega(r)}=-k_BT\log\rho(r)\Omega=k_BT\frac{r^2}{2r_0^2}+\tecxt\text{const}
\end{align*}
Mit einer zusätzlichen Kraft $f=-U'(r)$ und $U=-rf$ lautet die freie Energie
\begin{align*}
F&=-rf+k_BT\frac{r^2}{2r_0^2}
\end{align*} 
Das Minimum der freien Energie liegt dann bei $r_\tecxt\text{min}=\frac{r_0^2f}{k_BT}$ mit $F(r_\tecxt\text{min})=-\frac{r_0^2f^2}{2k_BT}$. Die Kraft nimmt dann die lineare Form $f=-\frac{k_BT}{r_0^2}\cdot r$.
\section*{Aufgabe 5.3: Maxwell-Relationen und Suszeptibilitäten}
\begin{table}[h!]
\centering
\begin{tabular}{c|c|c}
$-S$&$U$&$V$\\ \hline $H$&&$F$\\ \hline $-p$&$G$&$T$\\
\end{tabular}
\caption{Guggenheim-Quadrat}
\end{table}
Führe die abkürzende Notation $\frac{\partial A}{\partial B}=:\partial_BA=:A_{,B}$ ein. Soll eine Größe $C$ beim Ableiten von $A$ nach $B$ konstant gehalten werden, so schreiben wir $\kla\left( {A_{,B}} \right)_{C}=:A_{,B;C}$. Dann gilt:
\begin{align*}
T_{,V;S}&=U_{,S;V,V;S}=U_{,V;S,S;V}=p_{,S;V}\\
p_{,T;V}&=-F_{,V;T,T;V}=-F_{,T;V,V;T}=S_{,V;T}\\
S_{,p;T}&=-G_{,T;p,p;T}=-G_{,p;T,T;p}=-V_{,T;p}
\end{align*}
Das Kompressionsmodul lautet
\begin{align*}
K_T&:=-Vp_{,V;T}=VF_{,V;T,V;T}=V\kla\left( {\frac{\partial^2 F}{\partial V^2}} \right)_{T}
\end{align*}
Die Kompressibilität ist
\begin{align*}
\beta_T&:=-\frac{1}{V}V_{,p;T}=-\frac{1}{V}G_{,p;T,p;T}=-\frac{1}{V}\kla\left( {\frac{\partial^2G}{\partial p^2}} \right)_{T}
\end{align*}
Die spezifische Wärmekapazität kann man schreiben als
\begin{align*}
C_V&:=TS_{,T;V}=-TF_{,T;V,T;V}=-T\kla\left( {\frac{\partial^2F}{\partial T^2}} \right)_{V}
\end{align*}


\end{document}